% ============================================================
%  MTH636 -- Game Theory | Lecture 28
%  Behavior vs. Mixed Strategies & SPE Existence
% ============================================================
\documentclass[12pt, a4paper]{article}

% ============================================================
%  MTH636 -- Game Theory
%  Notes Preamble — Color-Coded, TikZ-Enhanced
% ============================================================

% --- Core packages ---
\usepackage[T1]{fontenc}
\usepackage[utf8]{inputenc}
\usepackage{lmodern}
\usepackage[margin=2.5cm, top=2.8cm, bottom=2.8cm]{geometry}
\usepackage{microtype}

% --- Math ---
\usepackage{amsmath, amssymb, amsthm, mathtools}
\usepackage{bm}        % bold math
\usepackage{cancel}    % strike-through in math

% --- Color & Boxes ---
\usepackage[dvipsnames, svgnames, x11names, table]{xcolor}
\usepackage[most]{tcolorbox}
\tcbuselibrary{breakable, skins, theorems}

% --- TikZ & Diagrams ---
\usepackage{tikz}
\usepackage{tikz-qtree}
\usetikzlibrary{arrows.meta, positioning, shapes.geometric, fit, backgrounds,
                decorations.pathreplacing, calc, matrix, trees}
\usepackage{pgfplots}
\pgfplotsset{compat=1.18}

% --- Tables ---
\usepackage{booktabs}
\usepackage{multirow}
\usepackage{array}
\usepackage{tabularx}
\usepackage{makecell}

% --- Typography & Layout ---
\usepackage{enumitem}
\usepackage{parskip}
\usepackage{fancyhdr}
\usepackage{titlesec}
\usepackage{titletoc}
\usepackage{hyperref}
\usepackage{graphicx}
\usepackage{float}
\usepackage{caption}

% --- Color Palette ---
\definecolor{defblue}{HTML}{1A5276}      % Definition header
\definecolor{defbluebg}{HTML}{D6EAF8}    % Definition background
\definecolor{thmgreen}{HTML}{1E8449}     % Theorem header
\definecolor{thmbg}{HTML}{D5F5E3}        % Theorem background
\definecolor{exorange}{HTML}{B7770D}     % Example header
\definecolor{exbg}{HTML}{FEF9E7}         % Example background
\definecolor{warnred}{HTML}{922B21}      % Warning/Key Insight header
\definecolor{warnbg}{HTML}{FADBD8}       % Warning background
\definecolor{noticegray}{HTML}{2C3E50}   % Note header
\definecolor{noticebg}{HTML}{EBF5FB}     % Note background
\definecolor{taketawayyellow}{HTML}{7D6608} % Key Takeaway header
\definecolor{takeawaybg}{HTML}{FDFEFE}   % Takeaway background
\definecolor{accentpurple}{HTML}{6C3483} % Accent / proof
\definecolor{lightgray}{HTML}{F2F3F4}    % shading

% ============================================================
%  TCOLORBOX ENVIRONMENTS
% ============================================================

% -- Definition (Blue) --
\newtcolorbox{defbox}[2][]{
  enhanced, breakable,
  colback=defbluebg, colframe=defblue, coltitle=white,
  fonttitle=\bfseries\sffamily,
  title={Definition: #2},
  attach boxed title to top left={yshift=-2mm, xshift=4mm},
  boxed title style={colback=defblue, rounded corners},
  arc=4pt, boxrule=1.2pt,
  left=6pt, right=6pt, top=8pt, bottom=6pt,
  #1
}

% -- Theorem / Proposition (Green) --
\newtcolorbox{thmbox}[2][]{
  enhanced, breakable,
  colback=thmbg, colframe=thmgreen, coltitle=white,
  fonttitle=\bfseries\sffamily,
  title={Theorem: #2},
  attach boxed title to top left={yshift=-2mm, xshift=4mm},
  boxed title style={colback=thmgreen, rounded corners},
  arc=4pt, boxrule=1.2pt,
  left=6pt, right=6pt, top=8pt, bottom=6pt,
  #1
}

% -- Proposition (Green variant) --
\newtcolorbox{propbox}[2][]{
  enhanced, breakable,
  colback=thmbg, colframe=thmgreen!70!black, coltitle=white,
  fonttitle=\bfseries\sffamily,
  title={Proposition: #2},
  attach boxed title to top left={yshift=-2mm, xshift=4mm},
  boxed title style={colback=thmgreen!70!black, rounded corners},
  arc=4pt, boxrule=1.2pt,
  left=6pt, right=6pt, top=8pt, bottom=6pt,
  #1
}

% -- Example (Orange/Yellow) --
\newtcolorbox{exbox}[2][]{
  enhanced, breakable,
  colback=exbg, colframe=exorange, coltitle=white,
  fonttitle=\bfseries\sffamily,
  title={Example: #2},
  attach boxed title to top left={yshift=-2mm, xshift=4mm},
  boxed title style={colback=exorange, rounded corners},
  arc=4pt, boxrule=1.2pt,
  left=6pt, right=6pt, top=8pt, bottom=6pt,
  #1
}

% -- Key Insight / Warning (Red) --
\newtcolorbox{warnbox}[2][]{
  enhanced, breakable,
  colback=warnbg, colframe=warnred, coltitle=white,
  fonttitle=\bfseries\sffamily,
  title={Key Insight: #2},
  attach boxed title to top left={yshift=-2mm, xshift=4mm},
  boxed title style={colback=warnred, rounded corners},
  arc=4pt, boxrule=1.2pt,
  left=6pt, right=6pt, top=8pt, bottom=6pt,
  #1
}

% -- Note (Gray-Blue) --
\newtcolorbox{notebox}[1][]{
  enhanced, breakable,
  colback=noticebg, colframe=noticegray!60, coltitle=white,
  fonttitle=\bfseries\sffamily,
  title={Note},
  attach boxed title to top left={yshift=-2mm, xshift=4mm},
  boxed title style={colback=noticegray!60, rounded corners},
  arc=4pt, boxrule=0.8pt,
  left=6pt, right=6pt, top=8pt, bottom=6pt,
  #1
}

% -- Key Takeaway --
\newtcolorbox{takebox}[1][]{
  enhanced, breakable,
  colback=takeawaybg, colframe=taketawayyellow, coltitle=white,
  fonttitle=\bfseries\sffamily,
  title={Key Takeaway},
  attach boxed title to top left={yshift=-2mm, xshift=4mm},
  boxed title style={colback=taketawayyellow, rounded corners},
  arc=4pt, boxrule=1.2pt,
  left=6pt, right=6pt, top=8pt, bottom=6pt,
  #1
}

% -- Intuition --
\newtcolorbox{intubox}[1][]{
  enhanced, breakable,
  colback=accentpurple!8, colframe=accentpurple!60, coltitle=white,
  fonttitle=\bfseries\sffamily,
  title={Intuition},
  attach boxed title to top left={yshift=-2mm, xshift=4mm},
  boxed title style={colback=accentpurple!80, rounded corners},
  arc=4pt, boxrule=0.8pt,
  left=6pt, right=6pt, top=8pt, bottom=6pt,
  #1
}

% ============================================================
%  PAYOFF MATRIX HELPER
%  Usage: inside tabular, use \payoff{a}{b} for (a,b)
% ============================================================
\newcommand{\payoff}[2]{$(#1,\;#2)$}
\newcommand{\nash}[2]{\underline{$#1$},\underline{$#2$}}  % underline NE payoffs
\newcommand{\Nset}{\mathbb{N}}
\newcommand{\Rset}{\mathbb{R}}
\newcommand{\Sset}{\mathcal{S}}
\newcommand{\Ui}{u_i}

% ============================================================
%  SECTION FORMATTING
% ============================================================
\titleformat{\section}
  {\large\bfseries\sffamily\color{defblue}}
  {\thesection}{1em}{}[\titlerule]

\titleformat{\subsection}
  {\normalsize\bfseries\sffamily\color{thmgreen!80!black}}
  {\thesubsection}{1em}{}

\titleformat{\subsubsection}
  {\normalsize\itshape\sffamily\color{exorange!80!black}}
  {\thesubsubsection}{1em}{}

% ============================================================
%  HEADER / FOOTER
% ============================================================
\pagestyle{fancy}
\fancyhf{}
\fancyhead[L]{\sffamily\small\color{defblue} MTH636 -- Game Theory}
\fancyhead[R]{\sffamily\small\color{defblue} \leftmark}
\fancyfoot[C]{\sffamily\small\thepage}
\renewcommand{\headrulewidth}{0.4pt}

% ============================================================
%  HYPERREF SETUP
% ============================================================
\hypersetup{
  colorlinks=true,
  linkcolor=defblue,
  urlcolor=thmgreen,
  citecolor=accentpurple,
  bookmarksnumbered=true,
  pdftitle={MTH636 Game Theory Notes},
  pdfauthor={MTH636 Student Notes}
}

% ============================================================
%  MISC HELPERS
% ============================================================
\newcommand{\BR}{\text{BR}}        % Best Response
\newcommand{\NE}{\text{NE}}        % Nash Equilibrium
\newcommand{\SPNE}{\text{SPNE}}    % Subgame Perfect NE
\newcommand{\Dom}{\text{dom}}      % Dominance
\newcommand{\mixs}[1]{\sigma_{#1}} % Mixed strategy
\newcommand{\pures}[1]{s_{#1}}     % Pure strategy


% ============================================================
\begin{document}

% Title
\begin{center}
  {\LARGE\bfseries\sffamily\color{defblue} MTH636 -- Game Theory}\\[6pt]
  {\Large\sffamily Lecture 28: Behavior Strategies, Mixed Strategies \& SPE Existence}\\[4pt]
\end{center}
\vspace{0.5cm}
\hrule height 1.5pt
\vspace{0.3cm}

% ============================================================
\section{Outcome Equivalence of Behavior and Mixed Strategies}

\begin{thmbox}{Outcome Equivalence Result}
In an extensive form game, every behavior strategy is outcome equivalent to a mixed strategy.
\end{thmbox}

\begin{proof}
Consider player $i$ and a behavior strategy $b_i$. Suppose player $i$ has $m$ many information sets $h_1, \ldots, h_m$.

Define the mixed strategy $\sigma_i$ induced by $b_i$ as follows:
\[
  \sigma_i(s_i) \;=\; \prod_{j=1}^{m} \beta_i\!\bigl(s_i(h_j),\, h_j\bigr)
  \qquad \forall\, s_i \in S_i.
\]

\subsection*{Validity of the Induced Mixed Strategy}

We first check that $\sigma_i$ is indeed a valid (probability) mixed strategy.

\medskip
\noindent\textbf{Non-negativity:} Since $\beta_i(s_i(h_j), h_j) \geq 0$ for all $j = 1, \ldots, m$, we have
\[
  \sigma_i(s_i) \;\geq\; 0.
\]

\noindent\textbf{Sums to one:}
\begin{align*}
  \sum_{s_i \in S_i} \sigma_i(s_i)
  &= \sum_{s_i \in S_i} \prod_{j=1}^{m} \beta_i\!\bigl(s_i(h_j),\, h_j\bigr) \\[4pt]
  &= \prod_{j=1}^{m} \sum_{s_i \in S_i} \beta_i\!\bigl(s_i(h_j),\, h_j\bigr) \\[4pt]
  &= \prod_{j=1}^{m} (1) \qquad \text{(since $\beta_i$ is a behavior strategy)} \\[4pt]
  &= 1.
\end{align*}

Hence $\sigma_i$ is a valid mixed strategy.

\subsection*{Outcome Equivalence}

Fix a leaf node $z$ and suppose $h_1, \ldots, h_k$ are the information sets of player $i$ on the path from the root to $z$. Also assume $a_1, \ldots, a_k$ are the actions of player $i$ at those information sets.

\medskip
\noindent Under $b_i$, player $i$'s contribution to the probability of reaching $z$ is:
\[
  \prod_{j=1}^{k} \beta_i(a_j,\, h_j).
\]

\noindent Under $\sigma_i$, the contribution becomes:
\[
  \sum_{\substack{s_i \in S_i \\ s_i(h_j) = a_j \\ \forall\, j = 1,\ldots,k}} \sigma_i(s_i).
\]

Expanding using the definition of $\sigma_i$:
\begin{align*}
  &= \sum_{\substack{s_i \in S_i \\ s_i(h_j) = a_j \\ \forall\, j = 1,\ldots,k}}
     \prod_{\ell=1}^{m} \beta_i\!\bigl(s_i(h_\ell),\, h_\ell\bigr) \\[6pt]
  &= \sum_{\substack{s_i \in S_i \\ s_i(h_j) = a_j \\ \forall\, j = 1,\ldots,k}}
     \prod_{\ell=1}^{k} \beta_i\!\bigl(s_i(h_\ell),\, h_\ell\bigr)
     \prod_{\ell=k+1}^{m} \beta_i\!\bigl(s_i(h_\ell),\, h_\ell\bigr) \\[6pt]
  &= \sum_{\substack{s_i \in S_i \\ s_i(h_j) = a_j \\ \forall\, j = 1,\ldots,k}}
     \prod_{\ell=1}^{k} \beta_i(a_\ell;\, h_\ell)
     \prod_{\ell=k+1}^{m} \beta_i\!\bigl(s_i(h_\ell);\, h_\ell\bigr) \\[6pt]
  &= \prod_{\ell=1}^{k} \beta_i(a_\ell;\, h_\ell)
     \sum_{\substack{s_i \in S_i \\ s_i(h_j) = a_j \\ \forall\, j = 1,\ldots,k}}
     \prod_{\ell=k+1}^{m} \beta_i\!\bigl(s_i(h_\ell);\, h_\ell\bigr) \\[6pt]
  &= \prod_{\ell=1}^{k} \beta_i(a_\ell;\, h_\ell)
     \prod_{\ell=k+1}^{m} \sum_{s_i \in S_i} \beta_i\!\bigl(s_i(h_\ell);\, h_\ell\bigr) \\[6pt]
  &= \prod_{\ell=1}^{k} \beta_i(a_\ell;\, h_\ell) \cdot \prod_{\ell=k+1}^{m} (1) \\[6pt]
  &= \prod_{\ell=1}^{k} \beta_i(a_\ell;\, h_\ell) \\[4pt]
  &= \prod_{j=1}^{k} \beta_i(a_j,\, h_j).
\end{align*}

Therefore $\sigma_i$ and $\beta_i$ are outcome equivalent.

\medskip
For the other part (every mixed strategy is outcome equivalent to a behavior strategy), a proof can be found in Solan's book.
\end{proof}

% ============================================================
\section{Existence of SPE in Extensive Form Games with Imperfect Information}

\begin{thmbox}{Existence of SPE}
Every extensive form game with imperfect information has a Subgame Perfect Equilibrium (SPE).
\end{thmbox}

\begin{intubox}
The method is \textbf{backward induction over subgames}. This will always result in an SPE.
\end{intubox}

\subsection{Example: Sequential Game with Imperfect Information}

Consider the following extensive form game. Player I moves first and chooses $A$ or $C$. If $A$ is chosen, Player II (who has an imperfect information set grouping two nodes) chooses $\ell$ or $r$ (from each of her two decision nodes). If $C$ is chosen, a chance node assigns probability $\tfrac{1}{4}$ to one branch and $\tfrac{3}{4}$ to another. The terminal payoffs are:

\begin{itemize}
  \item Under branch $A$: if II plays $\ell\ell$: $(4,1)$; if II plays $\ell r$: $(4,1)$ (from the left node's $\ell$) and $(7,2)$ (from the right node's $r$); etc.
  \item Specifically the four outcomes under Player I choosing $A$ and Player II's choices $(\ell,\ell)$, $(\ell,r)$, $(r,\ell)$, $(r,r)$ at her two-node information set yield terminal payoffs:
    $(4,1)$, $(7,2)$, $(6,3)$, $(4,10)$.
  \item Under branch $C$: with prob $\tfrac{1}{4}$, outcome $(8,4)$; with prob $\tfrac{3}{4}$, outcome $(8,8)$.
\end{itemize}

\noindent Expected payoff from branch $C$:
\[
  \tfrac{1}{4}(8,4) + \tfrac{3}{4}(8,8) = (8,\, \tfrac{1}{4}\cdot 4 + \tfrac{3}{4}\cdot 8) = (8,7).
\]

When solving by backward induction, Player I compares the payoff achievable under $A$ (given Player II's best response in her information set) with the payoff $(8,7)$ from $C$.

\medskip
In the reduced tree (after computing the continuation value), Player I chooses $A$ (giving access to the information-set subgame) or $C$ (giving $(8,7)$ directly).

\subsection{Backward Induction and SPE: Battle of the Sexes Example}

Consider the Man/Woman (M/W) sequential Battle of the Sexes game:
\begin{itemize}
  \item Man (M) moves first: $F$ (Football) or $D$ (Dance).
  \item Woman (W) observes and moves: $F$ or $D$ at each of her nodes.
\end{itemize}

\noindent The terminal payoffs are (Man, Woman):
\[
  (2,1),\quad (0,0),\quad (0,0),\quad (1,2)
\]
at the four leaves (corresponding to $(F,F)$, $(F,D)$, $(D,F)$, $(D,D)$).

\medskip
\noindent Nash equilibria of this game include:
\begin{itemize}
  \item $\bigl(F,\,(F,F)\bigr)$ -- this is an \textbf{NE} (Man plays $F$, Woman plays $F$ regardless).
  \item $\bigl(F,\,(F,D)\bigr)$ -- this is an \textbf{SPE} (Woman plays $F$ after $F$ and $D$ after $D$, which is sequentially rational).
\end{itemize}

\begin{notebox}
$(F,(F,D))$ is the unique SPE found by backward induction. The NE $(F,(F,F))$ is not an SPE because after the Man plays $D$, the Woman playing $F$ is not a best response in that subgame (she would get $0$ instead of $2$ by playing $D$).
\end{notebox}

% ============================================================
\section{Imperfect Information Game with Chance Nodes -- Further SPE Analysis}

\noindent Returning to the imperfect information example, consider a more elaborate tree where after Player I plays $A$ or $C$:
\begin{itemize}
  \item After $A$: Player II has a single information set (imperfect information -- she cannot distinguish the two nodes). Her actions are $\ell$ or $r$ at each node. Terminal payoffs:
    $(4,3)$,\; $(7,2)$,\; $(6,3)$,\; $(4,10)$.
  \item After $C$: a chance node (prob $\tfrac{1}{4}$ left, prob $\tfrac{3}{4}$ right). 
    Left subtree: Player II moves with actions $L,R$ yielding $(0_0)$, $(8_4)$.
    Right subtree: Player I moves with actions $a,b$ yielding $(8_8)$, $\left(\tfrac{0}{8}\right)$.
\end{itemize}

Working by backward induction over the subgames yields an SPE in which $(Ca,\, rR)$ is an SPE.

\begin{notebox}
\textbf{Perfect Bayesian Equilibrium} is another refinement (not covered in this syllabus).
\end{notebox}

% ============================================================
\section{From Strategic Form to Extensive Form}

\begin{itemize}
  \item \textbf{Question:} Given a strategic form game $G$, can we obtain an equivalent extensive form game $\Gamma$?
\end{itemize}

Formally:
\[
  G \;\longrightarrow\; \Gamma \;\longrightarrow\; G \quad \text{(equivalence)}
\]

Let
\[
  G = \bigl\langle N,\, (S_i)_{i \in N},\, (u_i)_{i \in N} \bigr\rangle
\]
and assume that $S_i = \{s_i^1, \ldots, s_i^{m_i}\}$ for each $i$.

\end{document}
